
本附录介绍带边黎曼曲面上的实\BookRecordIndex{T0210}柯西–黎曼算子，并证明这种情形下的\BookRecordIndex{T0221}黎曼--罗赫定理。第一节讨论基本背景并陈述该定理；C.2 节证明\BookRecordIndex{T0210}柯西–黎曼算子是弗雷德霍姆的；C.3 节介绍\BookRecordIndex{T0062}边界马斯洛夫指标；C.4 节证明\BookRecordIndex{T0221}黎曼--罗赫定理。C.5 节借助 Earle--Eells~\srcite{95} 的一个思路，说明如何由\BookRecordIndex{T0221}黎曼--罗赫定理推得\BookRecordIndex{T0222}黎曼映射定理。最后两节相对次要：前者简述非光滑丛上的\BookRecordIndex{T0210}柯西–黎曼算子，后者用\BookRecordIndex{T0358}辛流形切丛上可选取的各种联络来推导 $D_u$ 的公式。其结论已在注~\ref{rem:3-1-3} 中提及；它们并非本书其余部分所必需，只是为完整起见而收入。

在撰写本附录时，我们从与 Joel Robbin 的讨论中获益良多；他还慷慨同意将其对\BookRecordIndex{T0062}边界马斯洛夫指标的精彩阐述作为本书 C.3 节。

\section{\BookRecordIndex{T0210}柯西–黎曼算子}\label{sec:C-1}

设 $\Sigma$ 是带边紧致黎曼曲面，$E\to\Sigma$ 是光滑复向量丛。记 $j:T\Sigma\to T\Sigma$ 为 $\Sigma$ 上的复结构，$J:E\to E$ 为 $E$ 上的复结构。记
\[
\begin{aligned}
\Omega^k(\Sigma)&:\quad \Sigma\text{ 上光滑复值 $k$-形式的空间},\\
\Omega^{p,q}(\Sigma)&:\quad\text{其中 $(p,q)$ 型形式的空间},\\
\Omega^k(\Sigma,E)&:\quad\Sigma\text{ 上光滑 $E$ 值 $k$-形式的空间},\\
\Omega^{p,q}(\Sigma,E)&:\quad\text{其中 $(p,q)$ 型形式的空间}.
\end{aligned}
\]
这些空间分别是复向量丛 $\Lambda^kT^*\Sigma\otimes\mathbb C$、$\Lambda^{p,q}T^*\Sigma$、$\Lambda^kT^*\Sigma\otimes E$ 和 $\Lambda^{p,q}T^*\Sigma\otimes_{\mathbb C}E$ 的截面空间，因而是复向量空间。$\Sigma$ 上的复结构 $j$ 通过 $(a+ib,\alpha)\mapsto a\alpha+bj^*\alpha$ 在 $\Omega^k(\Sigma)$ 和 $\Omega^k(\Sigma,E)$ 上给出 $\mathbb C^*$ 作用；直和分解
\[
  \Omega^k(\Sigma)=\bigoplus_{p+q=k}\Omega^{p,q}(\Sigma),\qquad
  \Omega^k(\Sigma,E)=\bigoplus_{p+q=k}\Omega^{p,q}(\Sigma,E)
\]
的各项是这一作用的等型分量。将 $j$ 换为 $-j$ 会交换 $\Omega^{p,q}$ 与 $\Omega^{q,p}$。

记 $\bar\partial:\Omega^0(\Sigma)\to\Omega^{0,1}(\Sigma)$ 为外微分 $d:\Omega^0(\Sigma)\to\Omega^1(\Sigma)$ 与投影 $\Omega^1(\Sigma)=\Omega^{1,0}(\Sigma)\oplus\Omega^{0,1}(\Sigma)\to\Omega^{0,1}(\Sigma)$ 的复合。丛 $E\to\Sigma$ 上的（复线性、光滑）\emph{\BookRecordIndex{T0210}柯西–黎曼算子}是一个 $\mathbb C$-线性算子
\[
  D:\Omega^0(\Sigma,E)\longrightarrow\Omega^{0,1}(\Sigma,E),
\]
它满足 Leibniz 法则
\[
  D(f\xi)=f(D\xi)+(\bar\partial f)\xi
\]
（$\xi\in\Omega^0(\Sigma,E)$，$f\in\Omega^0(\Sigma)$）。例如，按坐标逐分量作用的 $\bar\partial$ 是平凡丛上的\BookRecordIndex{T0210}柯西–黎曼算子。算子 $D$ 唯一地延拓为满足同一 Leibniz 法则的斜楔导子 $D:\Omega^{p,q}(\Sigma,E)\to\Omega^{p,q+1}(\Sigma,E)$（名称相同，输入不同）。

\begin{remark}\label{rem:C-1-1}
丛 $E\to\Sigma$ 上的\BookRecordIndex{T0288}全纯结构是一个（极大）局部平凡化图册，其转换映射 $\phi$ 全纯。这样的结构唯一确定一个局部与 $\bar\partial$ 一致的\BookRecordIndex{T0210}柯西–黎曼算子 $D$，因为 $\bar\partial\phi=0$。每个（复线性、光滑）\BookRecordIndex{T0210}柯西–黎曼算子都以这种方式产生，这是一个深刻的定理（参见 Kobayashi~\srcite{212}）。因此，$E$ 上的\BookRecordIndex{T0288}全纯结构与 $E$ 上的\BookRecordIndex{T0210}柯西–黎曼算子一一对应。$E$ 上的\BookRecordIndex{T0288}全纯结构还确定全空间 $E$ 上的可积复结构：全纯局部平凡化给出 $E$ 上转换映射全纯的坐标图。
\end{remark}

\begin{remark}\label{rem:C-1-2}
$E$ 上的\emph{埃尔米特结构}是 $E$ 上一个实内积 $\langle\cdot,\cdot\rangle$，使复结构 $J$ 正交。$E$ 上的\emph{\BookRecordIndex{T0052}埃尔米特联络}是 $\mathbb C$-线性算子 $\nabla:\Omega^0(\Sigma,E)\to\Omega^1(\Sigma,E)$，使对所有 $f\in\Omega^0(\Sigma)$ 及 $\xi,\xi_1,\xi_2\in\Omega^0(\Sigma,E)$ 都有
\[
  \nabla(f\xi)=f\nabla\xi+(df)\xi,\qquad
  d\langle\xi_1,\xi_2\rangle
  =\langle\nabla\xi_1,\xi_2\rangle+\langle\xi_1,\nabla\xi_2\rangle.
\]
每个\BookRecordIndex{T0052}埃尔米特联络 $\nabla$ 都确定 $E$ 上的\BookRecordIndex{T0210}柯西–黎曼算子 $\bar\partial^\nabla$，定义为
\[
  \bar\partial^\nabla\xi:=(\nabla\xi)^{0,1}
  =\frac12\bigl(\nabla\xi+J\nabla\xi\circ j\bigr),
  \qquad \xi\in\Omega^0(\Sigma,E).
\]
反过来，对每个\BookRecordIndex{T0210}柯西–黎曼算子 $D$，存在唯一的\BookRecordIndex{T0052}埃尔米特联络 $\nabla$ 使 $D=\bar\partial^\nabla$。（证明这一点！）
\end{remark}

\begin{remark}\label{rem:C-1-3}
在 $E$ 上固定埃尔米特结构、\BookRecordIndex{T0052}埃尔米特联络 $\nabla$，并在 $\Sigma$ 上固定一个与 $j$ 确定的定向\BookRecordIndex{T0348}相容的体积形式 $\mathrm{dvol}$（即对 $0\ne v\in T\Sigma$，$\mathrm{dvol}(v,jv)>0$）。体积形式与埃尔米特结构在各空间 $\Omega^k(\Sigma,E)$ 上确定 $L^2$ 内积。我们断言，相对于这些内积，$\bar\partial^\nabla$ 的\BookRecordIndex{T0368}形式伴随算子为
\begin{equation}\label{eq:C-1-1}
  (\bar\partial^\nabla)^*=-*\partial^\nabla*=-*J\bar\partial^\nabla:
  \Omega^{0,1}(\Sigma,E)\longrightarrow\Omega^0(\Sigma,E).
\end{equation}
这里 $*$ 表示 \BookRecordIndex{T0026}Hodge 星算子；第二个等式来自对 $\eta\in\Omega^{0,1}(\Sigma,E)$ 有 $*\eta=-\eta\circ j=J\eta$。（若 $ds,dt$ 是 $T_z^*\Sigma$ 的正向标准正交基，则 $*ds=dt=-j^*ds$，$*dt=-ds=-j^*dt$。）\BookRecordIndex{T0368}形式伴随算子由恒等式
\begin{equation}\label{eq:C-1-2}
  \int_\Sigma\langle(\bar\partial^\nabla)^*\eta,\xi\rangle\,\mathrm{dvol}
  =\int_\Sigma\langle\eta,\bar\partial^\nabla\xi\rangle\,\mathrm{dvol}
\end{equation}
所刻画，其中 $\xi\in\Omega^0(\Sigma,E)$、$\eta\in\Omega^{0,1}(\Sigma,E)$，且一形式 $\langle\eta,J\xi\rangle$ 在 $T\partial\Sigma$ 上为零。因此，上述\BookRecordIndex{T0055}伴随公式由恒等式
\[
  -\langle*J\partial^\nabla\eta,\xi\rangle\,\mathrm{dvol}
  -\langle\eta,\bar\partial^\nabla\xi\rangle\,\mathrm{dvol}
  =d\langle\eta,J\xi\rangle
\]
得出。
\end{remark}

\begin{remark}\label{rem:C-1-4}
在 $\Sigma$ 上的局部全纯坐标中取 $E$ 的\BookRecordIndex{T0391}酉平凡化，记坐标为 $s+it\in\Omega\subset\mathbb H$，其中 $\mathbb H$ 是闭上半平面。在这种平凡化下，$J$ 等同于乘以 $i$，联络 $\nabla$ 写成
\[
  \nabla_s=\partial_s+\Phi,\qquad
  \nabla_t=\partial_t+\Psi,
\]
其中 $\Phi,\Psi:\Omega\to\mathfrak u(n)$。于是\BookRecordIndex{T0210}柯西–黎曼算子 $\bar\partial^\nabla$ 具有形式
\[
  \bar\partial^\nabla\xi
  =\frac12(\nabla_s\xi+i\nabla_t\xi)\,ds
   +\frac12(\nabla_t\xi-i\nabla_s\xi)\,dt,
  \qquad \xi:\Omega\to\mathbb C^n.
\]
$\Omega^{0,1}(\Omega,E)$ 中的元素写成 $\eta=\zeta\,ds-J\zeta\,dt$；对偶算子为
\[
  (\bar\partial^\nabla)^*(\zeta\,ds-i\zeta\,dt)
  =\lambda^{-2}(-\nabla_s\zeta+i\nabla_t\zeta),
\]
其中 $\Omega$ 上诱导的体积形式是 $\mathrm{dvol}=\lambda^2ds\wedge dt$。因此，把伴随算子视为作用于局部坐标中的函数 $\zeta$ 时，它就是 $\zeta\mapsto-2\lambda^{-2}\partial^\nabla\zeta$。
\end{remark}

\paragraph{实线性\BookRecordIndex{T0210}柯西–黎曼算子。} 在第 3 章的应用中，需要考虑更一般的\BookRecordIndex{T0210}柯西–黎曼算子，其零阶项是 $\mathbb R$-线性而非 $\mathbb C$-线性的。放宽光滑性假设也会有用。因此作如下定义。与注~\ref{rem:B-1-23} 一样，记 $W_F^{k,q}(\Sigma,E)$ 为
\[
  \Omega_F^0(\Sigma,E):=\{\xi\in\Omega^0(\Sigma,E)\mid\xi(\partial\Sigma)\subset F\}
\]
在\BookRecordIndex{T0319}索伯列夫空间 $W^{k,q}(\Sigma,E)$ 中的闭包；并记 $W_F^{k,q}(\Sigma,\Lambda^{0,1}T^*\Sigma\otimes E)$ 为
\[
  \Omega_F^{0,1}(\Sigma,E)
  :=\{\eta\in\Omega^{0,1}(\Sigma,E)\mid\eta(T\partial\Sigma)\subset F\}
\]
在 $W^{k,q}(\Sigma,\Lambda^{0,1}T^*\Sigma\otimes E)$ 中的闭包。

\begin{definition}\label{def:C-1-5}
设 $\ell$ 是正整数，$p>1$ 且 $\ell p>2$。$E$ 上一个 $W^{\ell-1,p}$ 类的\emph{实线性\BookRecordIndex{T0210}柯西–黎曼算子}是形如 $D=D_0+\alpha$ 的算子，其中 $\alpha\in W^{\ell-1,p}(\Sigma,\Lambda^{0,1}T^*\Sigma\otimes\operatorname{End}_{\mathbb R}(E))$，$D_0$ 是 $E$ 上光滑的复线性\BookRecordIndex{T0210}柯西–黎曼算子。
\end{definition}

实线性\BookRecordIndex{T0210}柯西–黎曼算子只对实值函数 $f$ 满足
\[
  D(f\xi)=f(D\xi)+(\bar\partial f)\xi.
\]
如复线性情形一样，它们也可由 $(E,J)$ 上的联络按公式
\[
  \bar\partial^\nabla\xi:=(\nabla\xi)^{0,1}
  =\frac12\bigl(\nabla\xi+J\nabla\xi\circ j\bigr)
\]
产生；但此时不再要求 $\nabla$ 是埃尔米特的，其\BookRecordIndex{T0223}联络势只需为 $W^{\ell-1,p}$ 类。因此，$\nabla$ 不必保持 $E$ 上的复结构或度量。更明确地说，任取 $E$ 上光滑的\BookRecordIndex{T0052}埃尔米特联络 $\nabla_0$，可写成
\[
  \nabla=\nabla_0+A,\qquad
  A\in W^{\ell-1,p}(\Sigma,T^*\Sigma\otimes_{\mathbb R}\operatorname{End}_{\mathbb R}(E)),
\]
从而 $D=\bar\partial^\nabla$ 具有形式
\[
  D\xi=D_0\xi+(A\xi)^{0,1},
\]
其中 $D_0=\bar\partial^{\nabla_0}$ 是与 $\nabla_0$ 关联的（复线性）\BookRecordIndex{T0210}柯西–黎曼算子。

\begin{exercise}\label{ex:C-1-6}
证明每个 $W^{\ell-1,p}$ 类的实线性\BookRecordIndex{T0210}柯西–黎曼算子，都等于某个 $W^{\ell-1,p}$ 类联络 $\nabla$ 的 $\bar\partial^\nabla$。
\end{exercise}

\begin{exercise}\label{ex:C-1-7}
证明每个实线性\BookRecordIndex{T0210}柯西–黎曼算子都\BookRecordIndex{T0213}可分解为一个复线性\BookRecordIndex{T0210}柯西–黎曼算子与一个零阶复反线性算子之和。
\end{exercise}

下面的习题表明，线丛上的每个复线性\BookRecordIndex{T0210}柯西–黎曼算子都与一个光滑的复线性\BookRecordIndex{T0210}柯西–黎曼算子\BookRecordIndex{T0153}规范等价。

\begin{exercise}\label{ex:C-1-8}
设 $E\to S$ 是闭黎曼曲面上的复线丛，$D$ 是 $S$ 上一个 $L^p$ 类复\BookRecordIndex{T0210}柯西–黎曼算子，$p>2$。证明
\[
  D=D_0+\alpha^{0,1},
\]
其中 $D_0$ 光滑且复线性，$\alpha\in L^p(S,T^*S\otimes_{\mathbb R}i\mathbb R)$。由 Hodge 理论，存在分解 $\alpha=\alpha_0+df+*dg$，其中 $f,g\in W^{1,p}(S,i\mathbb R)$，$\alpha_0\in\Omega^1(S,i\mathbb R)$ 为调和形式，因而光滑。证明由 $u:=e^{-f-ig}$ 给出的函数 $u\in W^{1,p}(S,\mathbb C^*)$ 满足 $u^{-1}\bar\partial u=-(df+*dg)^{0,1}$，从而
\[
  u^{-1}\circ D\circ u=D_0+\alpha_0^{0,1}.
\]
\end{exercise}

\paragraph{\BookRecordIndex{T0221}黎曼--罗赫定理。} 向量丛 $E\to\Sigma$ 上的黎曼--罗赫边值问题是算子
\[
  D_F:W_F^{\ell,p}(\Sigma,E)
  \longrightarrow W^{\ell-1,p}(\Sigma,\Lambda^{0,1}T^*\Sigma\otimes E),
\]
其中 $F\subset E|_{\partial\Sigma}$ 是全实子丛，而 $D_F$ 是 $W^{\ell-1,p}$ 类实线性\BookRecordIndex{T0210}柯西–黎曼算子 $D$ 在空间
\[
  W_F^{\ell,p}(\Sigma,E):=
  \{\xi\in W^{\ell,p}(\Sigma,E)\mid\xi(\partial\Sigma)\subset F\}
\]
上的限制。一如既往，假定 $p>1$ 且 $\ell p>2$。若 $E$ 上的埃尔米特内积 $\langle\cdot,\cdot\rangle$ 满足 $JF=F^\perp$，便称它与 $F$ \emph{\BookRecordIndex{T0348}相容}。相对于 $E$ 上这样的埃尔米特形式及 $\Sigma$ 上的体积形式 $\mathrm{dvol}$，$D_F$ 的对偶算子是\BookRecordIndex{T0368}形式伴随算子 $D^*$ 在子空间
\[
  W_F^{\ell,p}(\Sigma,\Lambda^{0,1}T^*\Sigma\otimes E)
  :=\{\eta\in W^{\ell,p}(\Sigma,\Lambda^{0,1}T^*\Sigma\otimes E)
        \mid\eta(T\partial\Sigma)\subset F\}
\]
上的限制，记作
\[
  D_F^*:W_F^{\ell,p}(\Sigma,\Lambda^{0,1}T^*\Sigma\otimes E)
  \longrightarrow W^{\ell-1,p}(\Sigma,E).
\]
以下引理是证明这些边界条件为弗雷德霍姆条件的关键；换言之，它们保证相应的算子 $D_F$ 与 $D_F^*$ 均为\BookRecordIndex{T0128}弗雷德霍姆算子（见定理~\ref{thm:C-2-3} 证明的步骤 3）。

\begin{lemma}\label{lem:C-1-9}
对任意 $\xi\in\Omega_F^0(\Sigma,E)$ 和 $\eta\in\Omega_F^{0,1}(\Sigma,E)$，
\[
  \int_\Sigma\langle D_F^*\eta,\xi\rangle\,\mathrm{dvol}
  =\int_\Sigma\langle\eta,D_F\xi\rangle\,\mathrm{dvol}.
\]
\end{lemma}
\begin{proof}
由于 $JF=F^\perp$，对所有 $\xi\in\Omega_F^0(\Sigma,E)$ 与 $\eta\in\Omega_F^{0,1}(\Sigma,E)$，一形式 $\langle\eta,J\xi\rangle$ 在 $T\partial\Sigma$ 上为零。因此，当 $D$ 光滑且复线性时，结论由注~\ref{rem:C-1-3} 得出。一般情形也直接成立，因为零阶算子 $\alpha$ 的形式伴随与其泛函分析意义下的伴随相同，即其转置 $\alpha^T$，逐点满足 $\langle\alpha^T\eta,\xi\rangle=\langle\eta,\alpha\xi\rangle$。
\end{proof}

现在陈述本附录的主定理。

\begin{theorem}[黎曼--罗赫]\label{thm:C-1-10}
设 $E\to\Sigma$ 是带边紧致黎曼曲面上的复向量丛，$F\subset E|_{\partial\Sigma}$ 是全实子丛。设 $D$ 是 $E$ 上 $W^{\ell-1,p}$ 类实线性\BookRecordIndex{T0210}柯西–黎曼算子，其中 $\ell$ 为正整数，$p>1$，且 $\ell p>2$。则对每个 $k\in\{1,\ldots,\ell\}$ 及每个满足 $k-2/q\le\ell-2/p$ 的实数 $q>1$，有：
\begin{enumerate}[label=(\roman*)]
\item 算子
\[
  D_F:W_F^{k,q}(\Sigma,E)\longrightarrow
  W^{k-1,q}(\Sigma,\Lambda^{0,1}T^*\Sigma\otimes E),
\]
\[
  D_F^*:W_F^{k,q}(\Sigma,\Lambda^{0,1}T^*\Sigma\otimes E)
  \longrightarrow W^{k-1,q}(\Sigma,E)
\]
是弗雷德霍姆的。此外，它们的核与 $k,q$ 无关，并有
\begin{equation}\label{eq:C-1-3}
  \eta\in\operatorname{im}D_F
  \quad\Longleftrightarrow\quad
  \int_\Sigma\langle\eta,\eta_0\rangle\,\mathrm{dvol}=0
  \quad\forall\eta_0\in\ker D_F^*
\end{equation}
（$\eta\in W^{k-1,q}(\Sigma,\Lambda^{0,1}T^*\Sigma\otimes E)$），以及
\begin{equation}\label{eq:C-1-4}
  \xi\in\operatorname{im}D_F^*
  \quad\Longleftrightarrow\quad
  \int_\Sigma\langle\xi,\xi_0\rangle\,\mathrm{dvol}=0
  \quad\forall\xi_0\in\ker D_F
\end{equation}
（$\xi\in W^{k-1,q}(\Sigma,E)$）。
\item $D_F$ 的实\BookRecordIndex{T0129}弗雷德霍姆指标为
\[
  \operatorname{index}(D_F)=n\chi(\Sigma)+\mu(E,F),
\]
其中 $\chi(\Sigma)$ 是 $\Sigma$ 的 Euler 示性数，$\mu(E,F)$ 是\BookRecordIndex{T0062}边界马斯洛夫指标（见 C.3 节），$n$ 是 $E$ 的复秩。当 $\partial\Sigma=\varnothing$ 时，\BookRecordIndex{T0062}边界马斯洛夫指标等于第一 Chern 数的两倍，即 $\mu(E)=2\langle c_1(E),[\Sigma]\rangle$。
\item 若 $n=1$，则
\[
  \mu(E,F)<0\quad\Longrightarrow\quad D_F\text{ 为单射},
\]
\[
  \mu(E,F)+2\chi(\Sigma)>0\quad\Longrightarrow\quad D_F\text{ 为满射}.
\]
\end{enumerate}
\end{theorem}

\paragraph{\BookRecordIndex{T0048}Serre 对偶。} 在定理~\ref{thm:C-1-10} 中，我们用 $\Sigma$ 上的黎曼度量和 $E$ 上的埃尔米特结构诱导的 $\Omega^0(\Sigma,E)$ 与 $\Omega^{0,1}(\Sigma,E)$ 的 $L^2$ 内积，把 $D_F:\Omega_F^0(\Sigma,E)\to\Omega^{0,1}(\Sigma,E)$ 的对偶算子写作从 $\Omega_F^{0,1}(\Sigma,E)$ 到 $\Omega^0(\Sigma,E)$ 的算子。这里给出复线性情形下对偶算子的另一种描述，它不要求 $\Sigma$ 上的体积形式或 $E$ 上的埃尔米特结构。考虑复向量丛
\[
  E^*:=\operatorname{Hom}^{\mathbb C}(E,\mathbb C)
\]
及其全实子丛
\[
  F^*:=\{(z,\zeta)\mid z\in\partial\Sigma,\
  \zeta\in\operatorname{Hom}^{\mathbb C}(E_z,\mathbb C),\
  \zeta(F_z)\subset\mathbb R\}\subset E^*|_{\partial\Sigma}.
\]
本小段中（且仅在本小段中），我们略微滥用记号，以 $T_z^*\Sigma=\operatorname{Hom}^{\mathbb C}(T_z\Sigma,\mathbb C)$ 表示 $T_z\Sigma$ 的复对偶空间，而不是本书其他地方的 $\Lambda^{1,0}T_z^*\Sigma$。因而
\[
  T^*\Sigma\otimes_{\mathbb C}E^*:=\operatorname{Hom}^{\mathbb C}(T\Sigma,E^*),
  \qquad
  T^*\partial\Sigma\otimes_{\mathbb R}F^*
  \subset\operatorname{Hom}^{\mathbb C}(T\Sigma,E^*)|_{\partial\Sigma}.
\]
$T^*\Sigma\otimes_{\mathbb C}E^*$ 的一个截面是 $\Sigma$ 上的 $E^*$ 值复线性一形式，即 $\Omega^0(\Sigma,T^*\Sigma\otimes_{\mathbb C}E^*)=\Omega^{1,0}(\Sigma,E^*)$；边界上全实子丛 $T^*\partial\Sigma\otimes_{\mathbb R}F^*$ 的截面，将边界切向量送入 $F^*$。存在自然的复双线性配对
\begin{equation}\label{eq:C-1-5}
  \Omega^0(\Sigma,T^*\Sigma\otimes_{\mathbb C}E^*)
  \times\Omega^{0,1}(\Sigma,E)
  \longrightarrow\mathbb C:
  (\zeta,\eta)\longmapsto\int_\Sigma\zeta\wedge\eta,
\end{equation}
其中二形式 $\zeta\wedge\eta\in\Omega^{1,1}(\Sigma)$ 由 $E^*$ 与 $E$ 的配对定义。依 Riesz 表示定理，此配对将 $T^*\Sigma\otimes_{\mathbb C}E^*$ 的 $L^2$ 截面空间等同于 $\Omega^{0,1}(\Sigma,E)$ 上对 $L^2$ 拓扑连续的复线性泛函空间。同样，有自然同构
\[
  \Omega^{0,1}(\Sigma,T^*\Sigma\otimes_{\mathbb C}E^*)
  \cong\Omega^{1,1}(\Sigma,E^*).
\]
具体而言，若 $\phi:T_z\Sigma\to\operatorname{Hom}^{\mathbb C}(T_z\Sigma,E_z^*)$ 是复反线性同态，对应的 $E_z^*$ 值二形式为
\begin{equation}\label{eq:C-1-6}
  \omega(\hat z_1,\hat z_2):=
  \phi(\hat z_1)(\hat z_2)-\phi(\hat z_2)(\hat z_1)
  \in E_z^*,\qquad \hat z_1,\hat z_2\in T_z\Sigma.
\end{equation}
逆映射将 $\omega\in\Lambda^{1,1}T_z^*\Sigma\otimes E_z^*$ 送到复反线性同态 $T_z\Sigma\to\operatorname{Hom}^{\mathbb C}(T_z\Sigma,E_z^*)$，即 $\hat z\mapsto\phi(\hat z):=\frac12\bigl(\omega(\hat z,\cdot)-i\omega(\hat z,i\cdot)\bigr)$。（习题：证明这些结论，并写出其局部坐标表达式。）因此又有复双线性配对
\begin{equation}\label{eq:C-1-7}
  \Omega^{0,1}(\Sigma,T^*\Sigma\otimes_{\mathbb C}E^*)
  \times\Omega^0(\Sigma,E)
  \longrightarrow\mathbb C:
  (\omega,\xi)\longmapsto\int_\Sigma\omega\wedge\xi,
\end{equation}
其中二形式 $\omega\wedge\xi\in\Omega^{1,1}(\Sigma)$ 由 $E^*$ 与 $E$ 的配对定义。

在 $E$ 上固定\BookRecordIndex{T0288}全纯结构，即取转换映射 $g_{\beta\alpha}:U_\alpha\cap U_\beta\to\mathrm{GL}(n,\mathbb C)$ 全纯的局部平凡化，并记相应\BookRecordIndex{T0210}柯西–黎曼算子为 $\bar\partial:\Omega^0(\Sigma,E)\to\Omega^{0,1}(\Sigma,E)$。$E^*$ 上有转换映射 $g_{\beta\alpha}^{T}$\sourcefootnote{原书将对偶丛的转换映射写为 $g_{\beta\alpha}^{T}$。若 $g_{\beta\alpha}$ 表示沿同一 $\alpha\to\beta$ 方向的纤维坐标变换，通常应为 $(g_{\beta\alpha}^{-1})^{T}$；原文未说明指标方向的约定，故此处保留原式而提示核对。} 的对应平凡化，它们确定\BookRecordIndex{T0210}柯西–黎曼算子 $\bar\partial:\Omega^0(\Sigma,E^*)\to\Omega^{0,1}(\Sigma,E^*)$。结合 $T^*\Sigma$ 上的\BookRecordIndex{T0288}全纯结构，得到\BookRecordIndex{T0210}柯西–黎曼算子
\[
  \bar\partial:\Omega^0(\Sigma,T^*\Sigma\otimes_{\mathbb C}E^*)
  \longrightarrow\Omega^{0,1}(\Sigma,T^*\Sigma\otimes_{\mathbb C}E^*).
\]
在 $\Omega^0(\Sigma,T^*\Sigma\otimes_{\mathbb C}E^*)\cong\Omega^{1,0}(\Sigma,E^*)$ 及 $\Omega^{0,1}(\Sigma,T^*\Sigma\otimes_{\mathbb C}E^*)\cong\Omega^{1,1}(\Sigma,E^*)$ 下，这就是延拓算子 $\bar\partial:\Omega^{1,0}(\Sigma,E^*)\to\Omega^{1,1}(\Sigma,E^*)$。现设
\[
  D=\bar\partial+A:\Omega^0(\Sigma,E)\longrightarrow\Omega^{0,1}(\Sigma,E)
\]
为一般\BookRecordIndex{T0210}柯西–黎曼算子，其中 $A\in\Omega^{0,1}(\Sigma,\operatorname{End}_{\mathbb R}(E))$，并记
\[
  D^*:=-\bar\partial+A^*:
  \Omega^0(\Sigma,T^*\Sigma\otimes_{\mathbb C}E^*)
  \longrightarrow\Omega^{0,1}(\Sigma,T^*\Sigma\otimes_{\mathbb C}E^*).
\]
（这里的记号不要与式~\eqref{eq:C-1-1} 中的\BookRecordIndex{T0368}形式伴随算子混淆。）于是 $D,D^*$ 满足恒等式\sourcefootnote{原书式 (C.1.8) 将右端的标量二形式误标为 $\Omega^{1,1}(\Sigma,E^*)$；由 $E^*$ 与 $E$ 的配对可知应属 $\Omega^{1,1}(\Sigma)$，此点亦见第二版 2019-11-27 勘误。}
\begin{equation}\label{eq:C-1-8}
  \zeta\wedge(D\xi)-(D^*\zeta)\wedge\xi
  =d(\zeta\wedge\xi)\in\Omega^{1,1}(\Sigma),
\end{equation}
其中 $\xi\in\Omega^0(\Sigma,E)$、$\zeta\in\Omega^0(\Sigma,T^*\Sigma\otimes_{\mathbb C}E^*)\cong\Omega^{1,0}(\Sigma,E^*)$，而一形式 $\zeta\wedge\xi\in\Omega^{1,0}(\Sigma)$ 仍通过 $E^*$ 与 $E$ 的配对理解。（习题：先设 $A=0$ 证明式~\eqref{eq:C-1-8}。）当 $A$ 仅为 $W^{\ell,p}$ 类时，此式仍成立。积分式~\eqref{eq:C-1-8} 得
\begin{equation}\label{eq:C-1-9}
  \int_\Sigma\zeta\wedge(D\xi)
  -\int_\Sigma(D^*\zeta)\wedge\xi
  =\int_{\partial\Sigma}\zeta\wedge\xi,
\end{equation}
其中 $\xi\in\Omega^0(\Sigma,E)$，$\zeta\in\Omega^0(\Sigma,T^*\Sigma\otimes_{\mathbb C}E^*)\cong\Omega^{1,0}(\Sigma,E^*)$。注意，若 $\xi(\partial\Sigma)\subset F$ 且 $\zeta(\partial\Sigma)\subset T^*\partial\Sigma\otimes_{\mathbb R}F^*$，则该一形式在 $\partial\Sigma$ 上的限制取实值。

\begin{corollary}[\BookRecordIndex{T0048}Serre 对偶]\label{cor:C-1-11}
设 $E\to\Sigma$ 为带边紧致黎曼曲面上的复向量丛，$F\subset E|_{\partial\Sigma}$ 为全实子丛。设 $D$ 是 $E$ 上一个 $W^{\ell-1,p}$ 类实线性\BookRecordIndex{T0210}柯西–黎曼算子，其中 $\ell$ 为正整数、$p>1$，且 $\ell p>2$。取 $\zeta\in L^r(\Sigma,T^*\Sigma\otimes_{\mathbb C}E^*)$，$r>1$。则以下两项等价：
\begin{enumerate}[label=(\roman*)]
\item 对每个 $\xi\in\Omega_F^0(\Sigma,E)$，$\int_\Sigma\zeta\wedge D\xi\in\mathbb R$；
\item $\zeta$ 为 $W^{\ell,p}$ 类，$D^*\zeta=0$，且 $\zeta|_{\partial\Sigma}$ 是子丛 $T^*\partial\Sigma\otimes_{\mathbb R}F^*$ 的截面。
\end{enumerate}
\end{corollary}
\begin{proof}
由上述讨论，(ii) 蕴涵 (i) 显然。为证明反向蕴涵，在 $\Sigma$ 上固定与复结构\BookRecordIndex{T0348}相容的黎曼度量，并如注~\ref{rem:C-1-2} 选取 $E$ 上使 $JF$ 与 $F$ 正交的实值埃尔米特内积 $\langle\cdot,\cdot\rangle_E$。记
\[
  \langle\cdot,\cdot\rangle_E^{\mathbb C}
  :=\langle\cdot,\cdot\rangle_E+i\langle J\cdot,\cdot\rangle_E
\]
为相应的复值埃尔米特形式：它对第一个自变量复反线性，对第二个自变量复线性。它诱导复反线性同构
\[
  \Omega_F^{0,1}(\Sigma,E)
  \longrightarrow\Omega^0_{T^*\partial\Sigma\otimes_{\mathbb R}F^*}
  (\Sigma,T^*\Sigma\otimes_{\mathbb C}E^*):
  \eta\longmapsto\zeta=\langle\eta,\cdot\rangle_E^{\mathbb C}.
\]
若 $\eta\in L^r(\Sigma,\Lambda^{0,1}T^*\Sigma\otimes_{\mathbb C}E)$，$\zeta:=\langle\eta,\cdot\rangle_E^{\mathbb C}$ 是相应的 $T^*\Sigma\otimes_{\mathbb C}E^*$ 的 $L^r$ 截面，则对每个 $\xi\in\Omega_F^0(\Sigma,E)$，
\[
  \operatorname{Im}\int_\Sigma\zeta\wedge D\xi
  =\int_\Sigma\langle J\eta\wedge D\xi\rangle_E
  =-\int_\Sigma\langle\eta\wedge JD\xi\rangle_E
  =-\int_\Sigma\langle\eta,D\xi\rangle_E\,\mathrm{dvol}_\Sigma.
\]
这里用到了 $JD\xi=-(D\xi)\circ j=*D\xi$，其中 $*$ 是一形式上的 \BookRecordIndex{T0026}Hodge 星算子。因此，若 $\zeta$ 满足 (i)，并选取 $q>1$ 使 $1/q+1/r=1$，则 $\eta\in L^r=(L^q)^*$ 消去算子
\[
  D_F:W_F^{1,q}(\Sigma,E)
  \longrightarrow L^q(\Sigma,\Lambda^{0,1}T^*\Sigma\otimes E)
\]
的像。由定理~\ref{thm:C-1-10}(i)，此像等于算子\sourcefootnote{原书对任意 $r>1$ 取其共轭指数 $q$，随即援引定理 C.1.10(i)；但该定理对 $k=1$ 还要求 $1-2/q\le\ell-2/p$，在 $\ell=1$ 且 $q>p$ 时不满足。此处保留原证明路线；未据此断定一般 $r$ 的结论错误，但该步尚需额外\BookRecordIndex{T0408}正则性论证。}
\[
  D_F^*:W_F^{1,r}(\Sigma,\Lambda^{0,1}T^*\Sigma\otimes E)
  \longrightarrow L^r(\Sigma,E)
\]
之核的湮灭子的湮灭子。由于自反巴拿赫空间中闭子空间的双重湮灭子就是它自身，$\eta$ 属于 $D_F^*$ 的核。再由定理~\ref{thm:C-1-10}(i)，此核与索伯列夫完备化的选择无关，故 $\eta\in W_F^{\ell,p}(\Sigma,\Lambda^{0,1}T^*\Sigma\otimes E)$。于是 $\zeta$ 为 $W^{\ell,p}$ 类，满足 $D^*\zeta=0$ 及所述边界条件，得证。
\end{proof}

\begin{remark}\label{rem:C-1-12}
对闭黎曼曲面上的光滑复线性\BookRecordIndex{T0210}柯西–黎曼算子，推论~\ref{cor:C-1-11} 断言，$\bar\partial:\Omega^0(\Sigma,T^*\Sigma\otimes_{\mathbb C}E^*)\to\Omega^{0,1}(\Sigma,T^*\Sigma\otimes_{\mathbb C}E^*)$ 的核，同构于 $\bar\partial:\Omega^0(\Sigma,E)\to\Omega^{0,1}(\Sigma,E)$ 的余核。这是 \BookRecordIndex{T0048}Serre 对偶的一个特例。一般地，\BookRecordIndex{T0048}Serre 对偶定理断言：若 $M$ 是复维数为 $n$ 的闭复流形，$E\to M$ 是全纯向量丛，则复向量空间间存在同构
\[
  H^q(M,\Omega^p(E^*))\cong H^{n-q}(M,\Omega^{n-p}(E))^*.
\]
这些是\BookRecordIndex{T0072}层上同调群；$\Omega^p(E)$ 表示 $M$ 上 $E$ 值全纯 $p$-形式的芽的\BookRecordIndex{T0072}层。
\end{remark}
